\documentclass[11pt,a4paper]{article}
\usepackage[utf8]{inputenc}
\usepackage[T1]{fontenc}
\usepackage[margin=2.3cm]{geometry}
\usepackage{helvet}
\renewcommand{\familydefault}{\sfdefault}
\usepackage[hidelinks]{hyperref}
\setlength{\parindent}{0pt}
\setlength{\parskip}{0.8em}
\begin{document}
\begin{center}
{\Large\bfseries GA-EXP-02 - FORM}\\[0.3em]
{\large Coevaluación de exposiciones}\\[0.5em]
ISR-701 - Aplicaciones Distribuidas\\
SPA 2026-2027\\
Byron Daniel Loor Mendoza\\
26 de noviembre de 2026
\end{center}
\section{Grupo GGHL - Computación paralela con Apache Spark}
GRUPO DEMORO UN POCO EN COMENZAR SU PRESENTACIÓN, A SU VEZ LES FALTO ASIGNAR UN CO-ANFRITIÓN PARA LA ACEPTACIÓN RAPIDA DE CADA ESTUDIANTE Y DOCENTE
POR PARTE DEL COMPAÑERO GIRALDO EXPLICA DE MANERA CORRECTA EL USO PRACTICO DE APACHE SPARK PARA INDAGAR EL TRABAJO DE SIMULACIÓN EN SU PROYECTO RTV, DISTRIBUCION E INICIO DE SESION, Y DEJO CLARO CADA UNA DE LAS PARTES ASIGNADAS EN SU PRACTICA EN BASE A SU PROYECTO DE FIN DE CURSO

EL COMPAÑERO ROGER INDAGA SOBRE CAP (COORDINAR DATOS AL EJECUTAR EN PARALELO), SU EXPLICACIÓN BASADA EN CLAB Y EJEMPLOS PUNTUALES INDAGADOS AL PARALELISMO, FUERON EXPRESADOS DE FORMA CORRECTA AL MENCIONAR LOS TIEMPOS DEL ESTADOS, TANTO EN AUTORIA Y COPIAS

EL COMPAÑERO LOMBEIDA EXPLICA LA TOLERANCIA A FALLOS, QUE LO ESTA EXPLICANDO DE MANERA CORRIDA, QUE MUY NOTORIAMENTE SE COMPRENDE PERO COMO SE DEBERIA, YA QUE NO TOMA EL TIEMPO SUFICIENTE DE INDAGAR Y EJEMPLIFICAR A TRAVES DE CODIGO O ALGUNA EXPLICACIÓN REAL
\end{document}

